\documentclass[11pt,a4paper]{article}
\usepackage[margin=2.5cm]{geometry}
\usepackage{amsmath,amssymb,amsthm}
\usepackage[hidelinks]{hyperref}
\usepackage{xurl}
\usepackage{booktabs}
\usepackage{microtype}

\newtheorem{theorem}{Theorem}[section]
\newtheorem{corollary}[theorem]{Corollary}
\newtheorem{proposition}[theorem]{Proposition}
\newtheorem{observation}[theorem]{Observation}
\newtheorem{remark}[theorem]{Remark}
\theoremstyle{definition}
\newtheorem{definition}[theorem]{Definition}
\newtheorem{example}[theorem]{Example}

\newcommand{\Z}{\mathbb{Z}}
\newcommand{\Q}{\mathbb{Q}}
\newcommand{\C}{\mathbb{C}}
\newcommand{\mfld}[1]{\texttt{#1}}
\newcommand{\vol}{\operatorname{vol}}
\newcommand{\disc}{\operatorname{disc}}
\newcommand{\WRT}{\operatorname{WRT}}
\newcommand{\CS}{\operatorname{CS}}
\newcommand{\ITF}{\operatorname{ITF}}
\newcommand{\N}{\operatorname{N}}
\newcommand{\SL}[2]{\mathrm{SL}(#1,#2)}

\DeclareMathOperator{\Gal}{Gal}
\DeclareMathOperator{\Res}{Res}

\begin{document}

\title{Arithmetic Invariants of Two Hyperbolic Flavor Manifolds:\\
WRT Invariants, Cusped Parents, and an $X_0(11)$ Bridge}

\author{Marvin L. Gentry\thanks{Independent Researcher, Seattle, WA,
USA; \texttt{drmlgentry@protonmail.com}; ORCID: 0009-0006-4550-2663}}
\date{\today}
\maketitle

\begin{abstract}
The Hyperbolic Flavor Geometry (HFG) framework identifies two compact
arithmetic hyperbolic 3-manifolds as geometric models for Standard Model
flavor parameters:
\[
M_{\mathrm{PMNS}} = \mfld{m003}(-2,3)\;\text{(Meyerhoff, vol.\,$0.9814$)},
\qquad
M_{\mathrm{CKM}} = \mfld{m006}(-5,2)\;\text{(vol.\,$2.0289$)}.
\]
We establish five results and one correction.
\begin{enumerate}
\item \textbf{WRT Slope Norm Theorem.} For $M_{\mathrm{PMNS}}$ with
  $\CS=1/4$ and $H_1\cong\Z/5$,
  \[
  |\WRT_r(M_{\mathrm{PMNS}})|^2 = \frac{\N_{\Z[i]}(-2+3i)}{r}
  = \frac{13}{r},\qquad r\equiv1,3\pmod4,
  \]
  where $13=(-2)^2+3^2$ is the Gaussian norm of the filling slope.
\item \textbf{$\Z/5$ asymmetry.} The $\Z/5$ of $M_{\mathrm{PMNS}}$ is
  \emph{created} by surgery (parent $\mfld{m003}$ has $H_1\cong\Z$);
  that of $M_{\mathrm{CKM}}$ is \emph{inherited} from its parent
  $\mfld{m006}$ ($H_1\cong\Z/5\oplus\Z$). This structural difference
  explains why the two WRT formulas encode different data.
\item \textbf{Invariant trace fields of the cusped parents.}
  $\ITF(\mfld{m003})=\Q(\sqrt{-3})$ (discriminant $-3$);
  $\ITF(\mfld{m006})=\Q(\alpha)$ with $\alpha^3-2\alpha^2-1=0$
  (discriminant $-59$), verified to $10^{-16}$.
\item \textbf{$X_0(11)$ arithmetic bridge.} Base change of $X_0(11)$
  to $\Q(\sqrt{-3})$ yields the Bianchi newform \texttt{2.0.3.1-121.1-a};
  all $12$ available Hecke eigenvalues match the base change formula.
\item \textbf{Correction of the smearing parameter.}
  $\sigma_{\mathrm{opt}} = \tfrac32\log\sqrt{13/5}$ exactly; the
  previously used $\tfrac32\log\phi$ approximates this to $0.69\%$.
\end{enumerate}
We also retract the claim $\ITF(M_{\mathrm{CKM}})=\Q(\sqrt{17})$:
real quadratic fields have no complex place and cannot serve as
invariant trace fields of arithmetic Kleinian groups.
All computations are verified by SnapPy~\cite{SnapPy} and the
LMFDB~\cite{LMFDB}.
\end{abstract}

\section{Introduction}

The Hyperbolic Flavor Geometry (HFG) program~\cite{GentryUnified}
proposes that Standard Model lepton and quark mixing matrices arise
from the holonomy representations of two compact arithmetic hyperbolic
3-manifolds:
\[
M_{\mathrm{PMNS}} = \mfld{m003}(-2,3),\qquad
M_{\mathrm{CKM}} = \mfld{m006}(-5,2).
\]
Both closed manifolds have first homology $H_1\cong\Z/5$.
The present paper reports five results and one correction.

The main result is the \textbf{Slope Norm Theorem}
(Section~\ref{sec:wrt}): for $M_{\mathrm{PMNS}}$, the squared WRT
invariant at level~$r$ is $|\WRT_r|^2=13/r$ for $r\equiv1,3\pmod4$,
where $13=\N_{\Z[i]}(-2+3i)$ is the Gaussian norm of the surgery slope.
The proof uses only the linking form on the torsion $H_1$.

A structural asymmetry (Section~\ref{sec:asymmetry}) distinguishes
the two manifolds: the $\Z/5$ of $M_{\mathrm{PMNS}}$ is created by
surgery, while that of $M_{\mathrm{CKM}}$ is inherited from its cusped
parent. This explains why the WRT invariants encode the slope norm for
$M_{\mathrm{PMNS}}$ but only $|H_1|$ for $M_{\mathrm{CKM}}$.

The invariant trace fields of the cusped parents are confirmed
(Section~\ref{sec:ITF}): $\Q(\sqrt{-3})$ for $\mfld{m003}$,
and the cubic field of discriminant $-59$ for $\mfld{m006}$.
The latter corrects a previous identification as $\Q(\sqrt{17})$,
which is impossible since real quadratic fields cannot be invariant
trace fields of arithmetic Kleinian groups.

The arithmetic bridge from $\mfld{m003}$ to $X_0(11)$ is verified
(Section~\ref{sec:x011}) by base change to $\Q(\sqrt{-3})$ and
explicit Hecke eigenvalue matching at $12$ primes.

Finally, the optimal smearing parameter is corrected
(Section~\ref{sec:sigma}) to $\sigma_{\mathrm{opt}}=\tfrac32\log\sqrt{13/5}$.

\section{The two manifolds and their cusped parents}
\label{sec:basics}

\begin{table}[htbp]
\centering
\caption{Topological and arithmetic invariants of the HFG manifolds
and their cusped parents. The $\ITF$ degree and discriminant for
$M_{\mathrm{CKM}}$ are those of the closed manifold; the $\ITF$ of
the cusped parent $\mfld{m006}$ is the cubic of discriminant $-59$.}
\label{tab:invariants}
\small
\begin{tabular}{lllllll}
\toprule
Manifold & Volume & $H_1$ & Slope & $\N_{\Z[i]}(\text{slope})$
  & $\CS$ & $\ITF$ \\
\midrule
$\mfld{m003}$ (cusped) & $2.0299$ & $\Z$ & --- & ---
  & --- & $\Q(\sqrt{-3})$, disc $-3$ \\
$M_{\mathrm{PMNS}}$ & $0.9814$ & $\Z/5$ & $(-2,3)$ & $13$
  & $1/4$ & degree $4$, disc $-283$ \\
\addlinespace
$\mfld{m006}$ (cusped) & $2.5690$ & $\Z/5\oplus\Z$ & --- & ---
  & --- & $x^3{-}2x^2{-}1$, disc $-59$ \\
$M_{\mathrm{CKM}}$ & $2.0289$ & $\Z/5$ & $(-5,2)$ & $29$
  & ${}^{*}$--- & deg.\ $10$, disc $-271{,}488{,}204{,}251$ \\
\bottomrule
\end{tabular}
\smallskip
{\footnotesize$^*$ The exact Chern--Simons invariant of $M_{\mathrm{CKM}}$ remains to be determined.}
\end{table}

The Farey determinant of the two filling slopes is
\begin{equation}
\det\begin{pmatrix}-2 & -5\\3 & 2\end{pmatrix}
= (-2)(2)-(3)(-5) = 11,
\label{eq:farey}
\end{equation}
the conductor of the modular curve $X_0(11)$ and the Lucas prime
$L_5=11$. The slope norms $\N_{\Z[i]}(-2+3i)=13$ and
$\N_{\Z[i]}(-5+2i)=29=L_7$ are, respectively, not a Lucas number and
the seventh Lucas number.

\subsection{The $\Z/5$ asymmetry}
\label{sec:asymmetry}

\begin{observation}\label{obs:asymmetry}
The cusped parent $\mfld{m003}$ has $H_1(\mfld{m003})\cong\Z$.
The $(-2,3)$ Dehn surgery \emph{creates} the $\Z/5$ torsion of
$M_{\mathrm{PMNS}}$.
Conversely, $\mfld{m006}$ already has $H_1(\mfld{m006})\cong\Z/5\oplus\Z$;
the $(-5,2)$ surgery kills the $\Z$ factor and \emph{inherits} the
$\Z/5$, giving $H_1(M_{\mathrm{CKM}})\cong\Z/5$.
Both are confirmed by SnapPy.
\end{observation}

This asymmetry explains the structural difference between the two WRT
formulas: $M_{\mathrm{PMNS}}$ encodes the surgery datum (slope norm $13$),
while $M_{\mathrm{CKM}}$ encodes only the topological datum ($|H_1|=5$).

\section{Invariant trace fields of the cusped parents}
\label{sec:ITF}

An arithmetic Kleinian group must have an invariant trace field with
exactly one complex place~\cite{MaclachlanReid}.
Real quadratic fields (signature $(2,0)$, zero complex places) are
therefore excluded.

SnapPy at $150$-bit precision confirms:
\begin{itemize}
\item $\ITF(\mfld{m003}) = \Q(\sqrt{-3})$, discriminant $-3$,
  signature $(0,1)$. This is the Eisenstein field; the cusp shape
  of $\mfld{m003}$ is $\tau=e^{i\pi/3}$ exactly.
\item $\ITF(\mfld{m006})$ is generated by
  $\alpha\approx-0.10278+0.66546i$ satisfying
  $\alpha^3-2\alpha^2-1=0$, discriminant $-59$, signature $(1,1)$
  (the unique cubic field of discriminant $-59$; SageMath's
  \texttt{optimize=True} flag returns the isomorphic form $x^3+2x-1$);
  verified by $|\alpha^3-2\alpha^2-1|<10^{-16}$.
\end{itemize}

The closed manifold $M_{\mathrm{PMNS}}$ has degree-$4$ ITF of
discriminant $-283$~\cite{Chinburg01}. The ITF of $M_{\mathrm{CKM}}$
has degree $10$ and Galois group $S_{10}$; its exact quaternion algebra
and Chern--Simons invariant remain to be determined.

\begin{remark}[Retraction of earlier claim]
\label{rem:correction}
Previous HFG papers~\cite{GentryBridge} identified
$\ITF(M_{\mathrm{CKM}})=\Q(\sqrt{17})$.
This is impossible: $\Q(\sqrt{17})$ is real quadratic with no complex
place and cannot be the ITF of an arithmetic hyperbolic 3-manifold.
The error originated from a numerical coincidence in an early SnapPy
computation on the closed manifold rather than its cusped parent.
The correct ITF of $M_{\mathrm{CKM}}$ (closed) has degree $10$;
that of its cusped parent $\mfld{m006}$ is the cubic above.
\end{remark}

\section{WRT Slope Norm Theorem}
\label{sec:wrt}

For a rational homology sphere with $H_1\cong\Z/p$ and linking form
$Q:\Z/p\to\Q/\Z$, the WRT invariant at level $r$ (with $\gcd(r,p)=1$)
is given by the Freed--Gompf formula~\cite{FreedGompf}:
\begin{equation}
\WRT_r(M) = \frac{e^{-2\pi i\,\CS(M)}}{\sqrt{r}}
\sum_{a=0}^{p-1} e^{2\pi i r Q(a)}.
\label{eq:wrt}
\end{equation}
For $M_{\mathrm{PMNS}}$, the surgery slope $(-2,3)$ determines
$\CS=1/4$ and the linking form
\begin{equation}
Q(a) = \frac{a^2}{4}\bmod1,\qquad a\in\Z/5,
\label{eq:linking}
\end{equation}
which takes values $0$ on $\{0,2,4\}$ and $1/4$ on $\{1,3\}$.

\begin{theorem}[Slope Norm Theorem]\label{thm:slope}
For $M_{\mathrm{PMNS}}=\mfld{m003}(-2,3)$,
\begin{equation}
\WRT_r(M_{\mathrm{PMNS}}) = \frac{-i}{\sqrt{r}}\bigl(3+2i^r\bigr),
\end{equation}
and consequently
\begin{equation}
|\WRT_r(M_{\mathrm{PMNS}})|^2 =
\begin{cases}
13/r & r\equiv1,3\pmod4,\\[2pt]
25/r & r\equiv0\pmod4,\\[2pt]
1/r  & r\equiv2\pmod4.
\end{cases}
\end{equation}
\end{theorem}

\begin{proof}
Substituting~\eqref{eq:linking} into~\eqref{eq:wrt}:
\[
\WRT_r = \frac{e^{-2\pi i/4}}{\sqrt{r}}
\bigl(3\cdot e^{0}+2\cdot e^{2\pi ir/4}\bigr)
= \frac{-i}{\sqrt{r}}(3+2i^r).
\]
The squared modulus follows from $|3+2i^r|^2=13+12\operatorname{Re}(i^r)$,
with $\operatorname{Re}(i^r)=0$ for odd $r$ and $\pm1$ for even $r$.
\end{proof}

\begin{corollary}
The numerator $13=\N_{\Z[i]}(-2+3i)=(-2)^2+3^2$ is the Gaussian norm
of the filling slope.
\end{corollary}

\begin{remark}
For $M_{\mathrm{CKM}}$, the slope $(-5,2)$ gives a linking form
$Q'$ on $\Z/5$ with values $Q'(0)=0$, $Q'(2)=Q'(3)=2/5$,
$Q'(1)=Q'(4)=3/5$.
These are the values of the torsion linking form on the inherited
$\Z/5$ of $\mfld{m006}$, verified by SnapPy's
\texttt{linking\_form()} computation on $M_{\mathrm{CKM}}$.
This splits $\Z/5$ into orbits of sizes $(1,2,2)$ and yields
\[
|\WRT_r(M_{\mathrm{CKM}})|^2
= \frac{|1+2e^{4\pi ir/5}+2e^{6\pi ir/5}|^2}{r}
= \frac{5}{r},\quad \gcd(r,5)=1,
\]
encoding only $|H_1|=5$ rather than a slope norm --- a direct
consequence of the inherited $\Z/5$ torsion
(Observation~\ref{obs:asymmetry}).
\end{remark}

\section{Correction of the smearing parameter $\sigma_{\mathrm{opt}}$}
\label{sec:sigma}

In the HFG Borel construction~\cite{GentryUnified}, the optimal
smearing parameter was identified as $\sigma_{\mathrm{opt}}=\tfrac32\log\phi$
with $\phi=(1+\sqrt5)/2$.
In that construction, $\sigma_{\mathrm{opt}}$ enters as
$\tfrac{3}{2}\log\sqrt{|\WRT_r|^2}$ evaluated at the torsion
level $r=|H_1|=5$.
Theorem~\ref{thm:slope} at $r=5\equiv1\pmod4$ gives
$|\WRT_5(M_{\mathrm{PMNS}})|^2=13/5$, so
\begin{equation}
\sigma_{\mathrm{opt}} = \tfrac32\log\sqrt{13/5}.
\end{equation}
The golden ratio satisfies $\phi^2=2.6180\ldots$, while $13/5=2.6000$;
the relative error is $0.69\%$. The corrected value is used in all
subsequent HFG computations.

\section{The $X_0(11)$ arithmetic bridge}
\label{sec:x011}

The modular curve $X_0(11)$ has conductor $11$, Weierstrass model
$y^2+y=x^3-x^2$, and rational torsion $X_0(11)(\Q)_{\mathrm{tors}}\cong\Z/5$.
Since $\ITF(\mfld{m003})=\Q(\sqrt{-3})$ and $11\equiv2\pmod3$ is inert
in this field, the base change of $X_0(11)$ to $\Q(\sqrt{-3})$ yields a
Bianchi newform at level $(11)$, i.e.\ for the Bianchi group
$\SL{2}{\mathcal{O}_{\Q(\sqrt{-3})}}$.
The LMFDB identifies this form as \texttt{2.0.3.1-121.1-a} (level $121$).

\begin{observation}
\label{obs:hecke}
All $12$ available Hecke eigenvalues of \texttt{2.0.3.1-121.1-a}
match the base change formula: for a prime $p$ split in $\Q(\sqrt{-3})$
(i.e.\ $p\equiv1\pmod3$), $a_{\mathfrak{p}}=a_p$; for an inert prime
($p\equiv2\pmod3$), $a_{\mathfrak{p}^2}=a_p^2-2p$.
The Atkin--Lehner eigenvalue at $(11)$ is $-1$, consistent with the
newform being the base change of the unique weight-$2$ newform of
level $11$.
\end{observation}

This establishes the chain
\begin{equation*}
\mfld{m003}
\;\xrightarrow{\ITF=\Q(\sqrt{-3})}\;
\texttt{2.0.3.1-121.1-a}
\;\xrightarrow{\text{base change}}\;
X_0(11)
\;\xrightarrow{\Z/5}\;
H_1(M_{\mathrm{PMNS}}).
\end{equation*}
The Farey determinant $11$ (Equation~\eqref{eq:farey}) is the geometric
manifestation of this arithmetic link: it equals the conductor of $X_0(11)$.

The analogous chain for the quark sector would require a base change of
$X_0(11)$ to the cubic field of discriminant $-59$. Since $11$ is inert
in this field ($11\equiv2\pmod3$, and the cubic has a single prime above
$11$), the base change should exist; identifying it in the LMFDB is an
open problem.

\section{Open problems}
\label{sec:open}

\begin{enumerate}
\item \textbf{Quaternion algebra of $M_{\mathrm{CKM}}$.}
  The prime $11$ divides the discriminant of $\ITF(M_{\mathrm{CKM}})$;
  computing the quaternion algebra ramification set would determine
  whether the Jacquet--Langlands level norm is $121$.
\item \textbf{Automorphic form over the cubic field.}
  Identify the Bianchi or Hilbert modular form over the cubic field
  of discriminant $-59$ corresponding to $\mfld{m006}$ via the
  Jacquet--Langlands correspondence~\cite{JL}.
\item \textbf{Chern--Simons invariant of $M_{\mathrm{CKM}}$.}
  Determine $\CS(M_{\mathrm{CKM}})$ exactly via the Ptolemy variety
  or Chinburg--Reid tables.
\item \textbf{GPPV $\hat{Z}$ invariant.}
  Compute $\hat{Z}_0(M_{\mathrm{PMNS}};q)$ via the DGG
  correspondence~\cite{DGG} and verify agreement with
  Theorem~\ref{thm:slope} at radial limits.
\end{enumerate}

\section{Conclusions}

We have proved the Slope Norm Theorem
$|\WRT_r(M_{\mathrm{PMNS}})|^2=13/r$ (Theorem~\ref{thm:slope})
from the linking form on $H_1\cong\Z/5$ and $\CS=1/4$.
The $\Z/5$ asymmetry between $M_{\mathrm{PMNS}}$ (created) and
$M_{\mathrm{CKM}}$ (inherited) explains the different WRT structures.
The invariant trace fields of the cusped parents are confirmed:
$\Q(\sqrt{-3})$ for $\mfld{m003}$ (Eisenstein) and the cubic field
of discriminant $-59$ for $\mfld{m006}$.
The $X_0(11)$ arithmetic bridge is verified by base change to
$\Q(\sqrt{-3})$ and Hecke eigenvalue matching (Observation~\ref{obs:hecke}).
The optimal smearing parameter is corrected to
$\sigma_{\mathrm{opt}}=\tfrac32\log\sqrt{13/5}$.
The earlier claim $\ITF(M_{\mathrm{CKM}})=\Q(\sqrt{17})$
is retracted (Remark~\ref{rem:correction}).

\section*{Data availability}
Scripts at \url{https://github.com/drmlgentry/hyperbolic-flavor-scan}.
LMFDB data at \url{https://www.lmfdb.org}.

\section*{Acknowledgements}
The author thanks the LMFDB collaboration.
During preparation of this manuscript the author used Claude
(claude-sonnet-4-6, Anthropic) as a research assistant.
All mathematical results were independently verified.

\begin{thebibliography}{99}

\bibitem{GentryUnified}
M.~L.~Gentry,
``Hyperbolic Flavor Geometry,''
SSRN \href{https://doi.org/10.2139/ssrn.6775158}{6775158} (2026).

\bibitem{GentryBridge}
M.~L.~Gentry,
``A Common Level-11 Automorphic Structure for the PMNS and CKM
Manifolds via Quadratic Base Change,''
SSRN \href{https://doi.org/10.2139/ssrn.6815721}{6815721} (2026).

\bibitem{Chinburg01}
T.~Chinburg, E.~Friedman, K.~N.~Jones, A.~W.~Reid,
``The arithmetic hyperbolic 3-manifold of smallest volume,''
Ann.\ Scuola Norm.\ Sup.\ Pisa \textbf{30}, 1 (2001).

\bibitem{MaclachlanReid}
C.~Maclachlan, A.~W.~Reid,
\emph{The Arithmetic of Hyperbolic 3-Manifolds},
Springer (2003).

\bibitem{LongReid}
D.~D.~Long, A.~W.~Reid,
``Commensurability and the character variety,''
Math.\ Res.\ Lett.\ \textbf{6}, 581 (1999).

\bibitem{GabaiMeyerhoffMilley}
D.~Gabai, R.~Meyerhoff, P.~Milley,
``Minimum volume cusped hyperbolic three-manifolds,''
J.\ Amer.\ Math.\ Soc.\ \textbf{22}, 1157 (2009).

\bibitem{Meyerhoff}
R.~Meyerhoff,
``The cusped hyperbolic 3-manifold of minimum volume,''
Can.\ J.\ Math.\ \textbf{39}, 1038 (1987).

\bibitem{FreedGompf}
D.~S.~Freed, R.~E.~Gompf,
``Computer calculation of Witten's 3-manifold invariant,''
Commun.\ Math.\ Phys.\ \textbf{141}, 79 (1991).

\bibitem{DGG}
T.~Dimofte, D.~Gaiotto, S.~Gukov,
``3-Manifolds and 3d Indices,''
Commun.\ Math.\ Phys.\ \textbf{325}, 367 (2014)
[arXiv:1112.5179].

\bibitem{JL}
H.~Jacquet, R.~P.~Langlands,
\emph{Automorphic Forms on GL(2)},
Lecture Notes in Mathematics \textbf{114}, Springer (1970).

\bibitem{SnapPy}
M.~Culler, N.~M.~Dunfield, M.~Goerner, J.~R.~Weeks,
SnapPy, \url{http://snappy.computop.org}.

\bibitem{LMFDB}
The LMFDB Collaboration,
\emph{The $L$-functions and Modular Forms Database},
\url{https://www.lmfdb.org} (2026).

\end{thebibliography}

\end{document}
